\subsection{沿抛物线的极大函数与 Hilbert 变换}
\label{sec:ch8-parabola-maximal-hilbert}

令 $L$ 为抛物型算子
\[
 Lu=\frac{\partial u}{\partial x_2}-\frac{\partial^2u}{\partial x_1^2}.
\]
作 Fourier 变换并仿照第~\ref{sec:ch4-operator-algebra}~节的论证, 得
\[
 \frac{\partial u}{\partial x_2}=T_1(Lu),
 \qquad
 \frac{\partial^2u}{\partial x_1^2}=T_2(Lu),
\]
其中
\[
 \widehat{T_1f}(\xi_1,\xi_2)
 =\frac{2\pi i\xi_2}{2\pi i\xi_2+4\pi^2|\xi_1|^2}\widehat f(\xi_1,\xi_2),
\]
\[
 \widehat{T_2f}(\xi_1,\xi_2)
 =\frac{4\pi^2|\xi_1|^2}{2\pi i\xi_2+4\pi^2|\xi_1|^2}\widehat f(\xi_1,\xi_2).
\]

\hypertarget{chapter-08-page-191}{}
这些乘子与第~\ref{sec:ch4-operator-algebra}~节所考虑的乘子不同, 并非零次齐次. 但它们满足
\[
 m(\lambda\xi_1,\lambda^2\xi_2)=m(\xi_1,\xi_2),
 \qquad \lambda>0.
\]
仿照命题~\ref{prop:ch4-fourier-homogeneity} 的证明可知, 若 $K(x_1,x_2)$ 是 $m$ 的 Fourier 变换, 则
\[
 K(\lambda x_1,\lambda^2x_2)=\lambda^{-3}K(x_1,x_2),
 \qquad \lambda>0.
\]

这个例子引导我们考虑以下类型的算子. 给定 $(x_1,x_2)\in\mathbb R^2\setminus\{(0,0)\}$, 存在唯一的 $r>0$ 和 $\theta\in S^1$ 使 $x_1=r\cos\theta$, $x_2=r^2\sin\theta$. 定义 $\Omega(\theta)=K(\cos\theta,\sin\theta)$, 则 $K(x_1,x_2)=r^{-3}\Omega(\theta)$. 定义
\[
 \begin{aligned}
 Tf(x_1,x_2)
 &=\operatorname{p.v.}\int_{\mathbb R^2}K(y_1,y_2)
 f(x_1-y_1,x_2-y_2)\,dy_1\,dy_2\\
 &=\operatorname{p.v.}\int_0^\infty\int_0^{2\pi}
 \frac{\Omega(\theta)}r
 f(x_1-r\cos\theta,x_2-r^2\sin\theta)
 (1+\sin^2\theta)\,dr\,d\theta.
 \end{aligned}
\]
若 $\Omega(\theta+\pi)=-\Omega(\theta)$, 即 $K$ 为奇函数, 则由旋转法中同样的论证, 见推论~\ref{cor:ch4-odd-kernel-lp},
\[
 Tf(x_1,x_2)=\int_0^\pi\Omega(\theta)H_\theta f(x_1,x_2)
 (1+\sin^2\theta)\,d\theta,
\]
其中
\[
 H_\theta f(x_1,x_2)=\operatorname{p.v.}\int_{\mathbb R}
 f(x_1-r\cos\theta,x_2-r^2\operatorname{sgn}(r)\sin\theta)\,\frac{dr}{r}.
\]
算子 $H_\theta$ 是沿奇抛物线
\[
 (r\cos\theta,r^2\operatorname{sgn}(r)\sin\theta),
 \qquad r\in\mathbb R,
\]
的 Hilbert 变换. 当 $\theta=0$ 或 $\pi/2$ 时, 这是一条直线, 因而得到通常的 Hilbert 变换. 由命题~\ref{prop:ch4-directional-operator-average} 的证明可知, 若 $\Omega\in L^1(S^1)$ 且算子 $H_\theta$ 在 $L^p$ 上一致有界, 则 $T$ 也有界.

为简单起见, 本节余下部分证明沿抛物线 $\Gamma(t)=(t,t^2)$ 的 Hilbert 变换与极大函数在 $L^p(\mathbb R^2)$, $1<p<\infty$, 上有界. 更确切地说, 考虑
\[
 H_\Gamma f(x_1,x_2)=\operatorname{p.v.}\int_{\mathbb R}
 f(x_1-t,x_2-t^2)\,\frac{dt}{t},
\]
\[
 M_\Gamma f(x_1,x_2)=\sup_{h>0}\left|\frac1{2h}
 \int_{-h}^hf(x_1-t,x_2-t^2)\,dt\right|.
\]
类似论证表明算子 $H_\theta$ 一致有界.

\hypertarget{chapter-08-page-192}{}
采用与上一节相似的方式分解 $H_\Gamma$ 和 $M_\Gamma$. 对 $j\in\mathbb Z$, 令有限测度 $\sigma_j$ 和 $\mu_j$ 对连续函数 $g$ 的作用为
\begin{equation}
 \sigma_j(g)=\int_{2^j<|t|<2^{j+1}}g(t,t^2)\,\frac{dt}{t},
 \label{eq:ch8-parabola-sigma-measure}
\end{equation}
\begin{equation}
 \mu_j(g)=\frac1{2^{j+1}}
 \int_{2^j<|t|<2^{j+1}}g(t,t^2)\,dt.
 \label{eq:ch8-parabola-mu-measure}
\end{equation}
于是
\begin{equation}
 H_\Gamma f=\sum_{j=-\infty}^{\infty}\sigma_j*f,
 \label{eq:ch8-parabola-hilbert-decomposition}
\end{equation}
并且
\begin{equation}
 M_\Gamma f\le2\sup_j\mu_j*|f|.
 \label{eq:ch8-parabola-maximal-decomposition}
\end{equation}
后一不等式成立, 因为若 $2^j\le h<2^{j+1}$, 则
\[
 \left|\frac1{2h}\int_{-h}^hf(x_1-t,x_2-t^2)\,dt\right|
 \le\frac1{2^{j+1}}\sum_{i=-\infty}^j2^{i+1}\mu_i*|f|.
\]

下面由两个更一般的结果推出 $H_\Gamma$ 和 $M_\Gamma$ 的有界性. 对测度 $\sigma$, 以 $|\sigma|$ 表示其全变差测度, 以 $\|\sigma\|$ 表示其全变差.

\begin{Theorem}
\label{thm:ch8-measure-sequence-operators}
设 $\{\sigma_j\}_{j\in\mathbb Z}$ 是有限 Borel 测度序列, $\|\sigma_j\|\le C$, 且对某个 $a>0$ 有
\begin{equation}
 |\widehat\sigma_j(\xi)|
 \le C\min(|2^j\xi_1|^a,|2^j\xi_1|^{-a}).
 \label{eq:ch8-measure-fourier-two-sided}
\end{equation}
若算子
\[
 \sigma^*(f)=\sup_j\bigl||\sigma_j|*f\bigr|
\]
在某个 $L^q$, $q>1$, 上有界, 则算子
\[
 Tf=\sum_{j=-\infty}^{\infty}\sigma_j*f,
 \qquad
 g(f)=\left(\sum_{j=-\infty}^{\infty}|\sigma_j*f|^2\right)^{1/2}
\]
在满足 $|1/p-1/2|<1/(2q)$ 的 $L^p$ 上有界.
\end{Theorem}

\begin{Proof}
定义算子 $S_j$ 为
\[
 \widehat{S_jf}(\xi)=\chi_{\Delta_j}(\xi_1)\widehat f(\xi),
 \qquad
 \Delta_j=(-2^{j+1},-2^j]\cup[2^j,2^{j+1}).
\]

\hypertarget{chapter-08-page-193}{}
则
\[
 Tf=\sum_{k=-\infty}^{\infty}T_kf,
 \qquad
 T_kf=\sum_{j=-\infty}^{\infty}\sigma_j*S_{j+k}f.
\]
仿照定理~\ref{thm:ch8-compact-kernel-dyadic-sum} 中证明~\eqref{eq:ch8-tk-l2-decay} 的论证, 由 Plancherel 定理和~\eqref{eq:ch8-measure-fourier-two-sided} 得
\begin{equation}
 \|T_kf\|_2\le C2^{-a|k|}\|f\|_2.
 \label{eq:ch8-measure-tk-l2}
\end{equation}
定义 $p_0$ 使 $1/2-1/p_0=1/(2q)$, 等价地 $q=(p_0/2)'$. 先证明
\begin{equation}
 \left\|\left(\sum_j|\sigma_j*g_j|^2\right)^{1/2}\right\|_{p_0}
 \le C\left\|\left(\sum_j|g_j|^2\right)^{1/2}\right\|_{p_0}.
 \label{eq:ch8-vector-measure-bound}
\end{equation}
因为 $|\sigma_j*g_j|^2\le\|\sigma_j\|(|\sigma_j|*|g_j|^2)$, 且~\eqref{eq:ch8-vector-measure-bound} 左端的平方是 $\sum_j|\sigma_j*g_j|^2$ 的 $p_0/2$ 范数, 由对偶性存在 $u\in L^q$, $\|u\|_q=1$, 使该范数等于
\[
 \int_{\mathbb R^2}\sum_j|\sigma_j*g_j|^2u
 \le C\sum_j\int_{\mathbb R^2}(|\sigma_j|*|g_j|^2)u
 \le C\sum_j\int_{\mathbb R^2}|g_j|^2\sigma^*(u).
\]
由 Hölder 不等式和关于 $\sigma^*$ 的假设得到~\eqref{eq:ch8-vector-measure-bound}. 因为 $S_iS_jf=0$ 除非 $i=j$, 由定理~\ref{thm:ch8-littlewood-paley} 和~\eqref{eq:ch8-vector-measure-bound},
\[
 \begin{aligned}
 \|T_kf\|_{p_0}
 &\le C\left\|\left(\sum_{j=-\infty}^{\infty}
 |\sigma_j*S_{j+k}f|^2\right)^{1/2}\right\|_{p_0}\\
 &\le C\left\|\left(\sum_{j=-\infty}^{\infty}|S_{j+k}f|^2\right)^{1/2}\right\|_{p_0}
 \le C\|f\|_{p_0}.
\end{aligned}
\]
在此式与~\eqref{eq:ch8-measure-tk-l2} 之间插值并对 $k$ 求和, 得 $T$ 在 $2\le p<p_0$ 上有界. 再由对偶性得到所需结论.

$g(f)$ 的有界性由与推论~\ref{cor:ch8-kernel-square-function} 的证明相同的论证得到.
\end{Proof}

\hypertarget{chapter-08-page-194}{}
\begin{Theorem}
\label{thm:ch8-parabolic-measure-maximal}
设 $\{\mu_j\}$ 是 $\mathbb R^2$ 上的正 Borel 测度序列, $\|\mu_j\|\le C$, 且对某个 $a>0$ 有
\begin{equation}
 |\widehat\mu_j(\xi)|\le C|2^j\xi_1|^{-a},
 \label{eq:ch8-mu-fourier-decay}
\end{equation}
\begin{equation}
 |\widehat\mu_j(\xi)-\widehat\mu_j(0,\xi_2)|
 \le C|2^j\xi_1|^a.
 \label{eq:ch8-mu-fourier-difference}
\end{equation}
令
\[
 \mathcal M_2f=\sup_j|\mu_j^2*f|
\]
为变量 $x_2$ 上的极大算子, 其中测度 $\mu_j^2$ 的 Fourier 变换为 $\widehat\mu_j(0,\xi_2)$. 若 $\mathcal M_2$ 在 $L^p(\mathbb R)$, $1<p\le\infty$, 上有界, 则极大算子
\[
 \mathcal Mf(x)=\sup_j|\mu_j*f(x)|
\]
也在 $L^p(\mathbb R^2)$, $1<p\le\infty$, 上有界.
\end{Theorem}

\begin{Proof}
固定 $\phi\in\mathcal S(\mathbb R)$, 使 $\widehat\phi(0)=1$, 并对每个 $j$ 定义测度 $\widetilde\sigma_j$ 为
\[
 \widehat{\widetilde\sigma_j}(\xi)
 =\widehat\mu_j(\xi)-\widehat\mu_j(0,\xi_2)\widehat\phi(2^j\xi_1).
\]
$\widetilde\sigma_j$ 满足定理~\ref{thm:ch8-measure-sequence-operators} 的假设. 对~\eqref{eq:ch8-measure-fourier-two-sided} 中负指数的估计, 使用~\eqref{eq:ch8-mu-fourier-decay} 并以 $Ct^{-a}$ 控制 $|\widehat\phi(t)|$. 对正指数的估计, 加上再减去 $\widehat\mu_j(0,\xi_2)$ 并应用中值定理. 令 $\widetilde g$ 和 $\widetilde\sigma^*$ 为相应算子. 仿照推论~\ref{cor:ch8-decaying-measure-maximal} 的证明, 得
\begin{equation}
 \mathcal Mf(x)\le\widetilde g(f)(x)+C\mathcal M_2M_1f(x),
 \label{eq:ch8-parabolic-maximal-bootstrap-one}
\end{equation}
\begin{equation}
 \widetilde\sigma^*(f)(x)\le\mathcal Mf(x)+C\mathcal M_2M_1f(x),
 \label{eq:ch8-parabolic-maximal-bootstrap-two}
\end{equation}
其中 $M_1$ 是第一个变量上的 Hardy--Littlewood 极大算子.

由 Plancherel 定理, $\widetilde g(f)$ 在 $L^2$ 上有界. 因此由~\eqref{eq:ch8-parabolic-maximal-bootstrap-one} 和关于 $\mathcal M_2$ 的假设, $\mathcal M$ 也在 $L^2$ 上有界. 再由~\eqref{eq:ch8-parabolic-maximal-bootstrap-two}, $\widetilde\sigma^*(f)$ 在 $L^2$ 上有界. 故由定理~\ref{thm:ch8-measure-sequence-operators}, $\widetilde g(f)$ 在 $4/3<p<4$ 上有界. 再用~\eqref{eq:ch8-parabolic-maximal-bootstrap-one} 和~\eqref{eq:ch8-parabolic-maximal-bootstrap-two}, $\mathcal M$ 与 $\widetilde\sigma^*$ 在同一范围内也有界. 定理~\ref{thm:ch8-measure-sequence-operators} 继而给出 $\widetilde g$ 在 $8/7<p<8$ 上有界. 反复应用这一提升论证, 得 $\mathcal M$, $\widetilde\sigma^*$ 和 $\widetilde g$ 在 $L^p$, $1<p<\infty$, 上有界.
\end{Proof}

为把这些定理用于 $H_\Gamma$ 和 $M_\Gamma$, 需要估计~\eqref{eq:ch8-parabola-hilbert-decomposition} 与~\eqref{eq:ch8-parabola-maximal-decomposition} 中测度 $\sigma_j$ 和 $\mu_j$ 的 Fourier 变换. 可用 Van der Corput 引理估计这些振荡积分. 因为这一结果本身也很重要, 下面给出证明.

\hypertarget{chapter-08-page-195}{}
\begin{Lemma}[Van der Corput's Lemma]
\label{lem:ch8-van-der-corput}
令
\[
 I(a,b)=\int_a^be^{ih(t)}\,dt.
\]
则:
\begin{enumerate}[label=(\arabic*)]
\item 若 $|h'(t)|\ge\lambda>0$ 且 $h'$ 单调, 则
\[
 |I(a,b)|\le C\lambda^{-1}.
\]
\item 若 $h\in C^k([a,b])$ 且 $|h^{(k)}(t)|\ge\lambda>0$, 则
\[
 |I(a,b)|\le C\lambda^{-1/k}.
\]
\end{enumerate}
两种情形中的常数均与 $a,b$ 无关.
\end{Lemma}

\begin{Proof}
(1) 分部积分得到
\[
 \begin{aligned}
 I(a,b)
 &=\int_a^bih'(t)e^{ih(t)}\frac{dt}{ih'(t)}\\
 &=\frac{e^{ih(b)}}{ih'(b)}-\frac{e^{ih(a)}}{ih'(a)}
 +i\int_a^be^{ih(t)}\,d\left(\frac1{h'(t)}\right).
 \end{aligned}
\]
因为 $h'$ 单调,
\[
 |I(a,b)|\le\frac2\lambda
 +\left|\int_a^b d\left(\frac1{h'(t)}\right)\right|
 \le\frac4\lambda.
\]

(2) 只证 $k=2$ 的情形, $k>2$ 时可由归纳法得到. 不妨设 $h''(t)\ge\lambda>0$. 则 $h'(t)$ 递增, 且至多在一个 $t_0\in(a,b)$ 处为零. 若这样的 $t_0$ 存在, 对 $\delta>0$ 定义
\[
 J_1=\{t\in(a,b):t<t_0-\delta\},
\]
\[
 J_2=(t_0-\delta,t_0+\delta)\cap(a,b),
 \qquad
 J_3=\{t\in(a,b):t>t_0+\delta\}.
\]
其中有些集合可能为空. 在 $J_1$ 和 $J_3$ 上, $|h'(t)|\ge\lambda\delta$, 故由 (1),
\[
 \left|\int_{J_1\cup J_3}e^{ih(t)}\,dt\right|
 \le8(\lambda\delta)^{-1}.
\]
$J_2$ 上的积分小于 $2\delta$. 取 $\delta=\lambda^{-1/2}$ 便得到所需结论. 若 $h'$ 在 $(a,b)$ 上从不为零, 则根据 $h'$ 为正还是为负, 以 $a$ 或 $b$ 代替 $t_0$ 作同样论证.
\end{Proof}

\begin{Lemma}
\label{lem:ch8-parabola-oscillatory-integral}
令
\[
 I(b)=\int_1^be^{i(t\xi_1+t^2\xi_2)}\,dt.
\]
则当 $1<b<2$ 且 $|\xi_1|>1$ 时,
\[
 |I(b)|\le C|\xi_1|^{-1/2}.
\]
\end{Lemma}

\hypertarget{chapter-08-page-196}{}
\begin{Proof}
令 $h(t)=t\xi_1+t^2\xi_2$. 当 $|\xi_1|\ge8|\xi_2|$ 时,
\[
 |h'(t)|=|\xi_1+2t\xi_2|\ge|\xi_1|/2.
\]
故由引理~\ref{lem:ch8-van-der-corput},
\[
 |I(b)|\le C|\xi_1|^{-1}\le C|\xi_1|^{-1/2},
\]
第二个不等式来自关于 $\xi_1$ 的假设. 另一方面, 因为 $h''(t)=2\xi_2$, 同一引理对任意 $\xi_1,\xi_2$ 给出
\[
 |I(b)|\le C|\xi_2|^{-1/2}.
\]
当 $|\xi_1|\le8|\xi_2|$ 时, 这蕴含 $|I(b)|\le C|\xi_1|^{-1/2}$.
\end{Proof}

现在可以用定理~\ref{thm:ch8-measure-sequence-operators} 和定理~\ref{thm:ch8-parabolic-measure-maximal} 证明所需结果.

\begin{Corollary}
\label{cor:ch8-parabola-hilbert-maximal}
算子 $H_\Gamma$ 和 $M_\Gamma$ 分别在 $L^p(\mathbb R^2)$, $1<p<\infty$, 和 $L^p(\mathbb R^2)$, $1<p\le\infty$, 上有界.
\end{Corollary}

\begin{Proof}
令 $\sigma_j,\mu_j$ 如~\eqref{eq:ch8-parabola-sigma-measure} 和~\eqref{eq:ch8-parabola-mu-measure} 中所定义. 由引理~\ref{lem:ch8-parabola-oscillatory-integral}, 当 $|\xi_1|>1$ 时,
\[
 |\widehat\mu_0(\xi)|,|\widehat\sigma_0(\xi)|
 \le C|\xi_1|^{-1/2}.
\]
其中 $\widehat\sigma_0$ 的估计还需使用分部积分. 再由关系
\[
 \widehat\mu_j(\xi)=\widehat\mu_0(2^j\xi_1,2^{2j}\xi_2),
 \qquad
 \widehat\sigma_j(\xi)=\widehat\sigma_0(2^j\xi_1,2^{2j}\xi_2),
\]
得到
\[
 |\widehat\mu_j(\xi)|,|\widehat\sigma_j(\xi)|
 \le C|2^j\xi_1|^{-1/2},
 \qquad |2^j\xi_1|>1.
\]
估计
\[
 |\widehat\sigma_j(\xi)|,
 |\widehat\mu_j(\xi)-\widehat\mu_j(0,\xi_2)|
 \le C|2^j\xi_1|
\]
立即由定义得到.

定理~\ref{thm:ch8-parabolic-measure-maximal} 中的极大算子 $\mathcal M_2$ 在此为
\[
 \sup_j\frac1{2^{j+1}}
 \left|\int_{2^j\le|t|\le2^{j+1}}g(x_2-t^2)\,dt\right|.
\]
作变量替换 $s=t^2$ 可知, $\mathcal M_2$ 由一维 Hardy--Littlewood 极大函数控制, 因而在 $L^p(\mathbb R)$, $1<p<\infty$, 上有界. 由定理~\ref{thm:ch8-parabolic-measure-maximal}, $M_\Gamma$ 在 $L^p(\mathbb R^2)$, $1<p\le\infty$, 上有界.

最后, 定理~\ref{thm:ch8-measure-sequence-operators} 中的极大算子 $\sigma^*$ 满足 $\sigma^*(f)\le M_\Gamma f$. 因而由该定理, $H_\Gamma$ 在 $L^p(\mathbb R^2)$, $1<p<\infty$, 上有界.
\end{Proof}
